
\medskip
\noindent \SGARefTarget{sga7.sga7i.eviii.n3-3}{3.3}.  We now turn to the case where $Q$ is finite locally
free over $S$, or is a group scheme of multiplicative type over $S$.
The key lemma in this case is the following well-known result.

\medskip
\noindent\underline{Proposition \SGARefTarget{sga7.sga7i.eviii.n3-3-1}{3.3.1}.}  \underline{Let $Q$ be a
commutative group scheme over $S$, and suppose that $Q$ is finite
locally free over $S$, or that $Q$ is of finite type and of
multiplicative type over $S$ (SGA 3 IX).  Then}
\[
 \uExt^1(Q,\underline G_m)=0.
\]

\noindent\underline{Proof.}  We distinguish the two cases.

\medskip
\noindent a) Suppose that $Q$ is finite locally free over $S$.  We
must prove that every extension $E$ of $Q$ by $\underline G_m$ splits
locally for the fppf topology.  We may plainly suppose $S$ affine;
there is then an integer $n>0$ such that ${}_nQ=0$ (SGA 3 IX 6.1).
Moreover, $n\cdot\id$ is an epimorphism of $\underline G_m$ for the
fppf topology, and hence the Kummer exact sequence
\[
 0\longrightarrow\mu_n\longrightarrow\underline G_m
 \xrightarrow{\ n\ }\underline G_m\longrightarrow0
\]
gives an exact sequence of the $\Ext^i(Q,-)$ and, in particular, an
\underline{epimorphism}
\[
 \Ext^1(Q,\mu_n)\longrightarrow
 {}_n\Ext^1(Q,\underline G_m)=\Ext^1(Q,\underline G_m).
\]
The last equality follows because multiplication by $n$, being zero
on $Q$, is also zero on $\Ext^1(Q,\underline G_m)$.  Thus the
extension $E$ comes from an extension $F$ of $Q$ by $\mu_n$.

The group $F$ is represented by a finite locally free scheme over $S$
(fpqc descent theory, SGA 1 VII).  By the known properties of Cartier
duality for finite locally free groups (SGA 3 VII 3),
\[
 D(F)=\uHom(F,\underline G_m)
\]
is an extension of $D(\mu_n)$ by $D(Q)$.  Here one uses the following
fact: if $A\longrightarrow B$ is a monomorphism of commutative finite
locally free groups over $S$, then
$D(B)\longrightarrow D(A)$ is faithfully flat.  Since $D(B)$ and
$D(A)$ are themselves finite locally free over $S$, the fibrewise
criterion for faithful flatness reduces us to the case where $S$ is
the spectrum of a field $k$.  In that case the assertion says that
the homomorphism of affine coordinate rings
\[
 k(D(A))\longrightarrow k(D(B))
\]
is injective (SGA 3 $\mathrm{VI}_B$ 11.14).  This homomorphism is the
transpose of $k(B)\longrightarrow k(A)$, which is surjective because
$A\longrightarrow B$ is a closed immersion (SGA 3 $\mathrm{VI}_B$
1.4.2).  Hence $k(D(A))\longrightarrow k(D(B))$ is indeed injective.

The inclusion $\mu_n\longrightarrow\underline G_m$ identifies with a
section of
\[
 D(\mu_n)=\uHom(\mu_n,\underline G_m),
\]
and it remains only to prove that, locally for the fppf topology on
$S$, this inclusion extends to a homomorphism
$F\longrightarrow\underline G_m$; equivalently, that the section of
$D(\mu_n)$ just considered lifts fppf-locally to a section of $D(F)$.
This is clear, since $D(F)\longrightarrow D(\mu_n)$ is faithfully flat
and of finite presentation.

\medskip
\noindent b) By SGA 3 X 4.5, after replacing $S$ by an \`etale cover,
we may suppose that $Q$ is diagonalizable, say $Q=D(M)$, where $M$ is
a finitely generated $\mathbb Z$-module.  Let $T$ be the torsion
submodule of $M$, and put $L=M/T$, so that $L\simeq\mathbb Z^r$.  The
exact sequence
\[
 0\longrightarrow T\longrightarrow M\longrightarrow L
 \longrightarrow0
\]
then gives an exact sequence (SGA 3 IX 3.1)
\[
 0\longrightarrow D(L)\longrightarrow D(M)\longrightarrow D(T)
 \longrightarrow0.
\]
It is therefore enough to prove separately that the $\uExt^1$ of
$D(L)$ and of $D(T)$ by $\underline G_m$ vanish.  The case of $D(T)$
follows from a), and the case of $D(L)$ reduces to $L=\mathbb Z$.
Thus it remains to check that
\[
 \uExt^1(\underline G_m,\underline G_m)=0,
\]
that is, that every extension $E$ of $\underline G_m$ by
$\underline G_m$ is fppf-locally trivial.

Such an extension is represented by an affine flat group scheme over
$S$ (fpqc descent theory, SGA 1 VII) whose fibres are tori; it is
therefore a torus (SGA 3 X 4.9).  By SGA 3 X 4.5 we may suppose that
$E$ is a split torus
\[
 E=D(R)\simeq\underline G_m^{,2}.
\]
After a Zariski localization, we may also suppose that the
monomorphism $\underline G_m\hookrightarrow E=D(R)$ comes from a
constant morphism $R\longrightarrow\mathbb Z_S$ (SGA 3 VIII 1.4),
which is necessarily an epimorphism (SGA 3 VIII 3.2 b)).  Choosing
$r\in R$ above $1\in\mathbb Z$ gives a homomorphism
$E\longrightarrow\underline G_m$ splitting the extension, and proves
the assertion.

\medskip
\noindent\underline{Corollary \SGARefTarget{sga7.sga7i.eviii.n3-3-2}{3.3.2}.}  \underline{Let $Q$ be as in
\SGARefLink{sga7.sga7i.eviii.n3-3-1}{3.3.1}.  Then, for every $P$, there is a canonical isomorphism}
\begin{equation}\tag{\SGARefTarget{sga7.sga7i.eviii.n3-3-2-1}{3.3.2.1}}
 \Biext^1(P,Q;\underline G_m)
 \xleftarrow{\ \sim\ }\Ext^1(P,D(Q)),
\end{equation}
\underline{where $D(Q)$ denotes the Cartier dual
$\uHom(Q,\underline G_m)$ of $Q$.  More precisely, the canonical
functor (\SGARefLink{sga7.sga7i.eviii.n1-1-6}{1.1.6}) is an equivalence from the category of extensions of
$P$ by $D(Q)$ to the category of biextensions of $(P,Q)$ by
$\underline G_m$.}

This is the special case of \SGARefLink{sga7.sga7i.eviii.n1-5}{1.5} obtained by applying \SGARefLink{sga7.sga7i.eviii.n3-3-1}{3.3.1}.

\medskip
\noindent\underline{Remarks \SGARefTarget{sga7.sga7i.eviii.n3-3-3}{3.3.3}.}  Under the hypotheses of \SGARefLink{sga7.sga7i.eviii.n3-3-1}{3.3.1},
$D(Q)$ is represented by a separated group scheme, locally of finite
presentation and locally quasi-finite over $S$.  More precisely, if
$Q$ is finite locally free over $S$, then so is $D(Q)$, as is well
known (SGA 3 VII 3).  If $Q$ is of multiplicative type, then $D(Q)$
is a \underline{finitely generated twisted constant group} (SGA 3 X
5.1), as was seen in SGA 3 IX 4.5 and X 5.7.  If $P$ is a group
scheme, every extension of $P$ by $D(Q)$ is therefore representable,
by a familiar lemma from fpqc descent theory (SGA 3 X 5.4).

\medskip
\noindent \SGARefTarget{sga7.sga7i.eviii.n3-3-4}{3.3.4}.  It is false that
$\uExt^1(Q,\underline G_m)=0$ for every affine flat group scheme $Q$
of finite presentation over $S$.  This already fails for
\[
 Q=\underline G_{a_S},\qquad
 S=\Spec k[t]/(t^2),
\]
where $k$ is a field of characteristic $p>0$.

